\documentclass[12pt,a4paper]{article}

\usepackage[utf8]{inputenc}
\usepackage[T1]{fontenc}
\usepackage[ngerman]{babel}
\usepackage{amsmath,amssymb,amsthm,mathtools}
\usepackage[colorlinks=true,linkcolor=blue,citecolor=red,urlcolor=blue]{hyperref}
\renewcommand*{\HyperDestNameFilter}[1]{\jobname-#1}
\usepackage{geometry}
\geometry{margin=2.5cm}
\usepackage{enumitem}
\usepackage{booktabs}
\usepackage{graphicx}
\usepackage{needspace}
\usepackage[capitalise,noabbrev]{cleveref}

\newtheorem{theorem}{Theorem}[section]
\newtheorem{lemma}[theorem]{Lemma}
\newtheorem{proposition}[theorem]{Proposition}
\newtheorem{corollary}[theorem]{Korollar}
\newtheorem{conjecture}[theorem]{Vermutung}
\theoremstyle{definition}
\newtheorem{definition}[theorem]{Definition}
\theoremstyle{remark}
\newtheorem{remark}[theorem]{Bemerkung}

\newcommand{\Ksym}{K_\sigma^{\mathrm{sym}}}
\newcommand{\EE}{\mathbb{E}}
\newcommand{\muG}{\mu_\sigma}
\newcommand{\PP}{\mathcal{P}}

\title{Die Riemannsche Vermutung: Conclusio}
\author{Lukas Geiger\\
\small Bernau, Deutschland\\
\small ORCID: \href{https://orcid.org/0009-0005-7296-1534}{0009-0005-7296-1534}}
\date{\today}

\begin{document}
\maketitle

\begin{abstract}\small
Diese abschließende Arbeit in einer dreiteiligen Serie führt die Ergebnisse der vorangegangenen Arbeiten zu einem einheitlichen Bild zusammen. Wir katalogisieren die bewiesenen strukturellen Ergebnisse (R1--R9), die ausgeschlossenen Wege (K1--K4) und die neuen Beiträge: das Shift-Parity-Lemma, die computergestützte finite-range Gap-Diagnostik für 33 Werte von $\lambda$ bis zu $1{,}300{,}000$, den Resolventen-gedämpften M1$''$-Rahmen, das Leading-Mode-Cancellation-Lemma mit der exakten Konstanten $c = 2 + \sqrt{2}$ und die Euler-Maclaurin-Proposition, die $\rho^{\mathrm{EM}}(L) < 0$ für $L \geq 14$ ($\lambda \geq 1{,}200{,}000$) etabliert. Die finite-range Diagnostik bleibt theoremstark bedingt auf eine validierte untere Schranke für den ungeraden Sektor. Teil~III berücksichtigt außerdem die in v2.0 eingeführten Nicht-Existenz-Theoreme (NE-A, NE-B in Teil~II), welche die vergleichende Analyse der Beweisstrategien neu strukturieren und bestätigen, dass sich gerade Dominanz nicht auf Totalpositivität oder Sturm--Liouville-artige Argumente reduzieren lässt. Wir präsentieren die Beweisarchitektur von Connes' Theorem 6.1 über die gerade Dominanz bis zur Riemannschen Vermutung als \emph{conditional reduction}, formulieren offene Probleme und diskutieren Verbindungen zur weiteren Landschaft der analytischen Zahlentheorie. In v2.1 wird ein robuster direkter Frontier-Dominanz-Beweis (Teil~II, Proposition~\emph{prop:direct-frontier}) hinzugefügt; in der v2.3-Lesart beweist er das asymptotische Variationsvorzeichen, aber nicht allein die Eigenwertordnung $\lambda_1^+<\lambda_1^-$.

\medskip\noindent
\textbf{v2.4 (23. Mai 2026): Parallelprogramm ergänzt.}  Ein neuer Unterabschnitt ``Parallelprogramm: Connes--vS-Fehlerkollaps und I-1-Normalfamilien-Reframe'' (\S\ref{sec:parallel-route-b}) präsentiert eine zweite Reduktion der Riemannschen Vermutung, unabhängig von der Asymptotischen Variationsgap-Vermutung~(A7), über die Connes--van-Suijlekom-err-Kollaps-Route (Route~B) und den I-1-Normalfamilien-Reframe der Begleitarbeit \cite{Geiger2026evendom}, \S E.1--E.18.  Beide Routen sind bedingt auf dieselbe offene Wand, Connes' \S6.6(b)/MS2-Streifenschranke~(ii), die der I-1-Reframe als (ii-a) + (ii-c) zerlegt und (modulo Einfachheit~(Z)) auf die einzige offene Schranke~(ii-a) absorbiert.  Konvergierte Empirie an sieben $\lambda \in \{3,5,7,9,11,13,15\}$ ist tabelliert, einschließlich der bit-identischen Zwei-dps-Validierung bei $\lambda = 11, 13, 15$.  Die DOI der Even-Dominance-Begleitarbeit wird vom v1.5-Record auf den v1.6-Record aktualisiert (\texttt{10.5281/zenodo.20291994}).
\medskip\noindent
\textbf{v2.5 (23. Mai 2026): Ausgeschlossene Wege des Spektralprogramms katalogisiert.}
Teil~III dokumentiert nun fünf nach v2.0 explizit ausgeschlossene Wege im
öffentlichen Audit-Trail: Eigenwertordnung über uniforme Spektrallücke, naive
Streifen-Lokalisierung, Rand-Profil-Faktorisierung, literales MS1 mit
Cluster-Aufspaltung sowie die zurückgezogenen Niedrigpräzisions-Messungen der
Streifen-Amplifikation. Die gemeinsame Obstruktionsstruktur ist der
nahentartete Minimal-Eigenvektor-Cluster von $QW_N$, weshalb tragfähige
Routen mit Cluster-/Quotienten-Reformulierungen und
wachstumsangepassten Normierungen arbeiten müssen statt mit
uniformen Gap- oder Dünn-Lokalisierungs-Heuristiken.
\end{abstract}

{\small
\noindent\textbf{MSC 2020:} 11M26, 11M36, 47A75, 65G20\\
\textbf{Schlagworte:} Riemannsche Vermutung, Li-Koeffizienten, gerade Dominanz, Euler-Maclaurin, Beweisarchitektur, offene Probleme
}

\clearpage
\tableofcontents

\clearpage

% =============================================================================
\section{Einleitung}
\label{sec:intro}
% =============================================================================

Diese Arbeit schließt eine dreiteilige Untersuchung der Riemannschen Vermutung durch Eichtheorie, Spektralanalyse und arithmetische Thermodynamik ab:
\begin{description}[font=\bfseries,leftmargin=2em]
  \item[Teil~I \cite{Geiger2026partI}:] Grundlagen und Hindernisse --- strukturelle Ergebnisse (R1--R9), vier kritische Hindernisse (K1--K4), Neuausrichtung auf das Spektralprogramm von Connes.
  \item[Teil~II \cite{Geiger2026partII}:] Gerade Dominanz --- Shift-Parity-Lemma, computergestützte finite-range Diagnostik ($33$ Werte, $\lambda$ bis zu $1{,}300{,}000$), Resolventen-gedämpfter M1$''$-Rahmen, Leading-Mode-Cancellation ($c = 2 + \sqrt{2}$), Euler-Maclaurin-Proposition ($\rho^{\mathrm{EM}} < 0$ für $\lambda \geq 1{,}200{,}000$).
  \item[Teil~III (diese Arbeit):] Conclusio --- Synthese, Beweisarchitektur, offene Probleme.
\end{description}

Die Kernbotschaft ist, dass die Riemannsche Vermutung, wenn sie über das Li-Kriterium und die quadratische Weil-Form angegangen wird, auf eine gut charakterisierte analytische Bedingung reduziert wird: eine asymptotische Variations-/Gap-Struktur des Kopplungsdifferenzials. Die Euler-Maclaurin-Proposition (Teil~II, Proposition~4.11) identifiziert das korrekte Variationsvorzeichen für großes $\lambda$; zusammen mit den $33$ computergestützten finite-range Diagnostiken (bis $\lambda = 1{,}300{,}000$) entsteht eine starke bedingte Brücke. Der noch offene Schritt ist die theoremstarke Übertragung dieser Variationsdiagnostik auf die Eigenwertordnung $\lambda_1^+<\lambda_1^-$.


% =============================================================================
\section{Katalog der bewiesenen Ergebnisse}
\label{sec:catalog}
% =============================================================================

\subsection{Strukturelle Ergebnisse (R1--R9)}

Die folgenden Ergebnisse sind unbedingt (sie setzen RH nicht voraus):
\begin{enumerate}[label=\textbf{R\arabic*},leftmargin=3em]
  \item \textbf{Gibbs-Rahmen.} Der Prime-Hub-Graph $G_\sigma$ mit dem Ruelle-Transferoperator $L_\sigma$ definiert ein wohldefiniertes Gibbs-Maß $\muG$ für alle $\sigma > 1$. Die Partitionsfunktion $P(\sigma) = -\log\zeta(\sigma)^{-1}$ divergiert für $\sigma \to 1^+$ (Teil~I, \S2).
  \item \textbf{Bombieri-Lagarias-Zerlegung.} $\lambda_n = \lambda_n^{(\infty)} + \lambda_n^{(\mathrm{fin})}$, wobei der archimedische Teil $\lambda_n^{(\infty)} \sim \frac{n}{2}\log n$ und der endliche Teil $|\lambda_n^{(\mathrm{fin})}| \leq C\sqrt{n}\log n$ unkonditioniert erfüllt (Teil~I, \S3).
  \item \textbf{Nicht-Identifikation.} Die Überschuss-Freie-Energie $I_n = \mathrm{KL}(\mu_\sigma^{(n)} \| \muG)$ erfüllt $I_n \geq 0$ unbedingt (als KL-Divergenz), aber $I_n \neq \lambda_n$. Die Gibbs-Normierung annihiliert den Beitrag des endlichen Teils (Teil~I, \S5).
  \item \textbf{Varianz-Schranke.} $\mathrm{Var}_{\muG}[X_{n,\sigma}] \leq \sum_p (\log p)^{2n}p^{-\sigma}$ für $\sigma > 1$. Dies beschränkt die Fluktuationen der thermodynamischen Li-Koeffizienten-Approximation (Teil~I, \S5).
  \item \textbf{Arithmetische Unschärferelation.} Für jede kompakt getragene Testfunktion $h$:
  \[ \mathrm{Var}[h(t)] \cdot \mathrm{Var}[\hat{h}(\omega)] \geq \frac{1}{4}\bigl|\langle h, h' \rangle\bigr|^2, \]
  dies verbindet die spektralen und räumlichen Unsicherheiten von Funktionen, die gegen die quadratische Weil-Form getestet werden (Teil~I, \S5).
  \item \textbf{Konvexität.} $\frac{d^2}{ds^2}\log\xi(\sigma) > 0$ nahe $\sigma = 1$. Die Funktion $\log\xi$ ist konvex, nicht konkav, was impliziert, dass $\lambda_1$ das Minimum (nicht das Maximum) des Ableitungsprofils ist. Dies widerlegt das ursprüngliche Monotonieargument (Teil~I, \S\S5 und~11).
  \item \textbf{OS-Reflexionspositivität.} Die Osterwalder-Schrader-Reflexion $\theta: t \mapsto -t$ erfüllt $\langle f, \theta f \rangle_{\muG} \geq 0$ für alle Testfunktionen $f$, die auf $t \geq 0$ getragen werden. Dies liefert eine Positivitätsstruktur im thermodynamischen Formalismus (Teil~I, \S5).
  \item \textbf{Spektrale Bestandsaufnahme.} Der arithmetische Kern $\Ksym$ auf $\ell^2(\PP_N)$ hat den negativen Trägheitsindex~$2$ für alle berechneten $N$ und $\sigma \in (1, 1.5)$. Genau zwei Eigenwerte sind negativ; die restlichen $N-2$ sind positiv (Teil~I, \S5).
  \item \textbf{Eich-Äquivalenz.} Für jede zulässige Eichung $g$ erfüllt der Kern $K_\sigma^g = \Ksym + \Delta K_g$ die Bedingung $\mathrm{sign}(\det K_\sigma^g) = \mathrm{sign}(\det \Ksym)$. Die Eichfreiheit ist tautologisch und kann die Vorzeichenstruktur nicht ändern (Teil~I, \S6).
\end{enumerate}


\subsection{Neue Ergebnisse aus den Teilen I--II}

\begin{enumerate}[label=\textbf{N\arabic*},leftmargin=3em]
  \item \textbf{Shift-Parity-Lemma} (Teil~II, \S2). Für alle $N \geq 3$ und alle $r \in (0,2)$:
  \[ \lambda_1(D_N(r)) < 0, \]
  wobei $D_N(r) = S_N^+(r) - S_N^-(r)$ die Differenz zwischen den Shift-Matrizen des geraden und ungeraden Sektors ist. Analytisch bewiesen (Zwei-Fall-det/Spur-Argument für $N=3$, Monotoniereduktion für $N \geq 4$) und durch Intervallarithmetik zertifiziert (55.000 Teilintervalle, 60-stellige Präzision, Gerschgorin-Kreisschranken).

  \item \textbf{Computergestützte finite-range Diagnostik} (Teil~II). Für $33$ Werte von $\lambda$ von $100$ bis $1{,}300{,}000$ ist eine starke negative Gap-Diagnostik durch Intervallarithmetik dokumentiert (mpmath.iv, 50-stellige Präzision). Die Lücke wächst monoton als $\sim\!\sqrt{\lambda}$, von $-2{,}25$ bei $\lambda = 100$ bis $-475{,}83$ bei $\lambda = 1{,}300{,}000$. Als theoremstarke gerade Dominanz bleibt dies bedingt auf die validierte untere Schranke für den ungeraden Sektor.

  \item \textbf{Universelle gerade Präferenz einzelner Primzahlen} (Teil~II, Korollar~2.5). Für jede Primzahl $p$ mit $r = \log p / \log\lambda \in (0,2)$:
  \[ \lambda_1(W_p^+) \leq \lambda_1(W_p^-) + \lambda_1(D_N(r)) < \lambda_1(W_p^-), \]
  nach dem Shift-Parity-Lemma und der Weyl-Ungleichung.

  \item \textbf{Frontier-Primzahl-Mechanismus} (Teil~II, \S2). Die gerade Dominanz wird von Primzahlen $p \sim \lambda$ (normierter Shift $r \approx 1$) getrieben, nicht von kleinen Primzahlen. Mit wachsendem $\lambda$ steigt die Anzahl der Frontier-Primzahlen als $\pi(\lambda) - \pi(\lambda/e) \sim \lambda/\log\lambda$.

  \item \textbf{Hellmann-Feynman-Analyse} (Teil~II, \S2). Die $L$-Ableitung $\partial\Delta/\partial L > 0$ für alle $\lambda \geq 80$ bis $\lambda = 500$. Die Vertiefung der Lücke wird durch neue Primzahlen getrieben, die in die Summe eintreten (Primzahl-Zählmechanismus), nicht durch das Wachstum von $L = \log\lambda$. Dies schließt $L$-Monotonie als Beweisstrategie aus.

  \item \textbf{Lücken-Asymptotik} (Teil~II). Die gerade-ungerade Lücke skaliert als $\Delta(\lambda) \sim -c\sqrt{\lambda}$ (exponentiell in $L = \log\lambda$), robust negativ für alle $\lambda \geq 55$, mit $|\Delta| \to \infty$.

  \item \textbf{Resolventen-gedämpfter M1$''$-Rahmen} (Teil~II). Das Kopplungsdifferenzial wird mittels Störungstheorie zweiter Ordnung zerlegt. Die Resolventen-gedämpfte Energiedifferenz erfüllt $(E_{\sin} - E_{\cos})/D(\lambda) = O(1/L) \to 0$, wodurch der verbleibende analytische Schritt auf Fourier-Analyse reduziert wird.

  \item \textbf{Leading-Mode-Cancellation-Lemma} (Teil~II). Die Überlapp-Differenzen zwischen geradem und ungeradem Sektor heben sich paarweise bei $u = 0$ auf (Fourier-Orthogonalität), wobei die Ableitung die exakte Konstante $c = 2 + \sqrt{2} \approx 3{,}414$ liefert. Die geschlossenen Profile $S^+_{\mathrm{sym}}(1,0;u) = \sqrt{2}\sin(\pi u)/\pi$ und $S^+_{\mathrm{sym}}(1,2;u) = \tfrac{2}{3\pi}(4\cos\pi u - 1)\sin\pi u$ werden symbolisch verifiziert.

  \item \textbf{Euler-Maclaurin-Proposition} (Teil~II, Proposition~4.11). Die Primsummen, die $B_j$ und $g_j$ definieren, werden durch ihre PNT-asymptotischen Integrale ersetzt, um das Euler-Maclaurin-Verhältnis $\rho^{\mathrm{EM}}(L)$ zu erhalten. Dann ist $\rho^{\mathrm{EM}}(L)$ monoton fallend für $L \geq 7$ und hat eine Nullstelle bei $L \approx 13{,}5$. Für $L \geq 14$ (d.\,h. $\lambda \geq e^{14} \approx 1{,}200{,}000$):
  \[
    \rho^{\mathrm{EM}}(L) < 0 \quad\text{(d.\,h. } E_{\sin}^{\mathrm{EM}} < E_{\cos}^{\mathrm{EM}}\text{)},
  \]
  so dass $\mathrm{cert\_gap}^{\mathrm{EM}}(\lambda) < 0$.

  \item \textbf{PNT-Transfer-Lemma und M1$''$-Reduktion} (Teil~II). Durch Abel-Summation mit der Chebyshev-Funktion $\theta(x) = x + O(x\,e^{-c\sqrt{\log x}})$ erfüllt das Resolventen-Verhältnis $\rho(\lambda) = \rho^{\mathrm{EM}}(\lambda) + O(L\,e^{-c\sqrt{L}})$. Da $\rho^{\mathrm{EM}} < 0$ für $L \geq 14$ und der Fehler verschwindet, liefert dies eine asymptotische Variationsreduktion des Resolventen-Subdominanz-Mechanismus. Der theoremstarke Restschritt bleibt die Odd-Sektor-Lower-Bound-Schließung, die das Variationsvorzeichen in die Eigenwertungleichung $\lambda_1^+ < \lambda_1^-$ überführt.
\end{enumerate}


% =============================================================================
\section{Katalog der ausgeschlossenen Wege}
\label{sec:excluded}
% =============================================================================

Vier unabhängige Hindernisse machen die ursprüngliche eichtheoretische Beweisstrategie zunichte:
\begin{enumerate}[label=\textbf{K\arabic*},leftmargin=3em]
  \item \textbf{Gibbs-Normierungslücke.} Die $1/P(\sigma)$-Normierung annihiliert den Beitrag des endlichen Teils $\lambda_1^{(\mathrm{fin})} = \gamma$. Der thermodynamische Li-Koeffizient $m_n(\sigma)$ konvergiert nicht gegen $\lambda_n$ (Teil~I, \S9).
  \item \textbf{Vorzeichenwechsel von $F(\sigma)$.} Das Zielfunktional $F(\sigma) = A(\sigma)B(\sigma) - P(\sigma)[C(\sigma) + D(\sigma)]$ wechselt bei $\sigma_0 \approx 1{,}14$ das Vorzeichen. Keine zulässige Eichung kann die Positivität über die Randschicht $(1, 1{,}14)$ hinaus wiederherstellen (Teil~I, \S10).
  \item \textbf{Falsche Monotonierichtung.} Die Konvexität von $\log\xi$ nahe $s = 1$ bedeutet, dass $\lambda_1$ das Minimum des Ableitungsprofils ist, nicht das Maximum. Das Monotonieargument läuft in die falsche Richtung (Teil~I, \S11).
  \item \textbf{Spurklassen-Hindernis.} Der Euler-Produkt-Operator $\mathcal{L}_s = \mathrm{diag}(p^{-s})$ ist nur für $\Re(s) > 1$ von der Spurklasse ($\sum_p p^{-1/2} = +\infty$). Die Fredholm-Determinante $\det(I - \mathcal{L}_s)$ existiert nicht auf der kritischen Geraden. Dies ist die fundamentale Dichtelücke zwischen Primzahlen und Geodäten (Teil~I, \S12).
\end{enumerate}

\begin{remark}[Die Hindernisse sind unabhängig]
Jedes der Hindernisse K1--K4 ist individuell fatal. K1 tötet die Gibbs-Identifikation, K2 tötet den Eich-Ansatz, K3 kehrt das Monotonieargument um, K4 tötet die spektrale Determinantenkonstruktion. Selbst wenn drei behoben würden, genügte das vierte, um die ursprüngliche Strategie zu blockieren.
\end{remark}

Zusätzlich wurden die folgenden Wege erkundet und als unzureichend befunden:
\begin{itemize}
  \item \textbf{de-Branges-Positivität:} Conrey und Li (2000) zeigten, dass die Kern-Positivitätsbedingung für die relevanten Hilbert-Räume, die mit $\zeta$ assoziiert sind, \emph{nicht} erfüllt ist.
  \item \textbf{Stieltjes-Realisierung:} Die Folge $c_n = \lambda_n/n$ ist streng wachsend, was die vollständige Monotonie verletzt. Das Spektralmaß unter RH lebt auf $|w| = 1$, nicht auf $[0,1]$.
  \item \textbf{$L$-Monotonie der Lücke:} Die Hellmann-Feynman-Analyse (N5) zeigt $\partial\Delta/\partial L > 0$ für $\lambda \geq 80$, was einen Beweis der Lückenvertiefung allein über $L$-Differentiation ausschließt.
  \item \textbf{Absolute totale Positivität ($\mathrm{PF}_\infty$) für $A_\lambda$ (v1.4):} Der Fourier-Multiplikator des Primzahl-Shift-Operators ist $M(\xi) = -2\,\Re[\zeta'/\zeta(\tfrac{1}{2}+i\xi)]$, eine Identität aus der expliziten Weil-Formel, gültig ohne RH-Annahme. Da $M$ auf ungefähr $49\%$ seiner Domäne negativ ist (die nicht-trivialen Nullstellen von $\zeta$ erzeugen jeweils eine Oszillation), ist der Primzahl-Shift-Operator nicht positiv-semidefinit, und $A_\lambda$ ist kein $\mathrm{PF}_\infty$. Schwache Variationsverminderung im primzahlgewichteten Mittel gilt und fällt mit dem Frontier-Dominanz-Argument (N4 + N6) zusammen, aber es gibt kein absolutes ,,Primzahlen glätten alles''-Prinzip.
  \item \textbf{Universeller kommutierender Operator (v1.4):} Eine SVD-Suche nach symmetrischem $T$ mit $[T, D_N(r)] = 0$ simultan auf einem dichten $r$-Gitter in $(0,2)$ zeigt für $N \in \{3,5,7,10,15\}$, dass der Kern eindimensional ist und ausschließlich aus Vielfachen der Identität besteht. Der zweitkleinste Singulärwert bleibt bei $N = 15$ unterhalb von $\sim\!0{,}70$ nach unten beschränkt und nimmt mit $N$ nicht ab. Die Primzahl-Shifts auf verschiedenen Skalen tragen daher keine gemeinsame verborgene Symmetrie. Dies schließt jedes Sturm-Liouville-artige Einfachheits-Argument aus und bestätigt, dass der Beweismechanismus arithmetisch-statistisch ist (primzahlgewichtete Summation der punktweisen Shift-Parity), nicht algebraisch.
\end{itemize}

% added 2026-05-23: spectral-programme excluded paths (v2.5 catalog extension)
\subsection*{Ausgeschlossene Wege im Spektralprogramm (v2.5-Katalog)}
\label{sec:excluded-spectral}

Das Post-v2.0-Spektralprogramm (Connes--van Suijlekom-Reformulierung, I-1 Normalfamilien-Reframe, err-Kollaps-Diagnostik) hat ebenfalls mehrere Wege geschlossen. Diese werden hier für den Audit-Trail dokumentiert; sie betreffen nicht die conditional reduction (die auf A7 ruht).

\begin{itemize}
  \item \textbf{Eigenwertordnung über uniforme Spektrallücke.} Der Plan, $\lambda_{\min}(W^+-W^-)<0$ über Kato--Temple-, Davis--Kahan- oder Lehmann-Gap-Abschätzungen auf die Eigenwertungleichung $\lambda_1^+<\lambda_1^-$ zu heben, scheitert daran, dass der relevante Minimal-Eigenvektor von $QW_N$ in einem nahezu entarteten Cluster sitzt (even/odd-Splitt $\sim\!10^{-65}$ bei $\lambda=5$, $\sim\!10^{-84}$ bei $\lambda=7$, super-exponentiell mit $\lambda$ kollabierend). Keine uniforme Spektrallücke verfügbar; gap-abhängige Methoden können die Schranke nicht liefern. Begleitend: EvenDom \texttt{K0\_CCM\_CLUSTER\_LADDER}.

  \item \textbf{Lokalisierungs-Hypothese für die Streifenschranke (ii-a).} In Route~B (Connes--vS err-Kollaps / I-1 Reframe) war der natürliche Plan, $|\hat\xi_\lambda(x+i\delta)|$ über die Lokalisierung $|\hat\xi_\lambda(z)|\leq\int|\xi_\lambda(t)|e^{|y||t|}\,dt$ zu schranken, unter der Annahme, dass $\xi_\lambda$ in einem $\lambda$-uniformen Kern nahe $t=0$ konzentriert ist. Falsifiziert: $\xi_\lambda$ ist genuin \emph{randkonzentriert} ($r_{99}\approx L/2$, $|\xi_\lambda(0)|\approx 0$ --- die $\delta_N$-Unsichtbarkeit geometrisch gesehen). Die naive (ii-a) ist daher nicht nur schwer, sondern in $L^2$-Normierung falsch; eine wachstumsangepasste Normierung $\Phi$ ist nötig. Begleitend: EvenDom \texttt{K0\_NORMAL\_FAMILY\_REFRAME} \S E.8 (\texttt{eigenfunction\_decay\_audit.py}).

  \item \textbf{Rand-Profil-Faktorisierung für (ii-a).} Die exakte Identität $F_\lambda(z) = 2\cos(zL/2)\,C_\lambda(z) + 2\sin(zL/2)\,S_\lambda(z)$ isoliert die explizite $\lambda^\delta$-Wachstumsmodulation in $\cos/\sin(zL/2)$ und würde (ii-a) auf eine uniforme Typ-Schranke an $\{C_\lambda,S_\lambda\}$ reduzieren --- vorausgesetzt, das Randprofil $\eta_\lambda(s)=\xi_\lambda(L/2-s)$ ist in einer dünnen $\lambda$-uniformen Schicht konzentriert. Empirisch ist die Eigenfunktionsmasse ausgebreitet ($r_{50}\approx 1$ uniform, aber $r_{99}\approx L/2$); die Edge-Route liefert keinen Typ-Gewinn. Begleitend: EvenDom \texttt{K0\_NORMAL\_FAMILY\_REFRAME} \S E.9.

  \item \textbf{Literales MS1 (Cluster-Leiter).} Die konservative Form von MS1 --- eine uniform einfache Aufspaltung des minimalen Eigenwerts von $QW_N$ an jedem $(\lambda,N)$ --- scheitert: Der Splitt kollabiert super-exponentiell mit $\lambda$. Ein literales MS1 ist nicht verfügbar; nur eine Cluster-/Quotienten-Reformulierung hat strukturelle Traktion. Begleitend: EvenDom \texttt{K0\_CCM\_MS1\_ROUTE\_S\_AUDIT}, \texttt{K0\_CCM\_CLUSTER\_MS1\_REFINEMENT}.

  \item \textbf{Streifen-Amplifikation bei $\mathrm{dps}=30$ (retrahiert).} Frühe Messungen der Streifen-Amplifikation $A_\lambda(\delta)=\sup_K|\hat\xi_\lambda(x+i\delta)|/\sup_K|\hat\xi_\lambda(x)|$ bei $\mathrm{dps}=30$ suggerierten sub-generisches Wachstum und nicht-monotone $N/R$-Instabilität (\S\S E.16--E.17). Bei konvergierter Präzision ($\mathrm{dps}\approx 30\text{--}50\cdot\lambda$, nötig wegen der Cluster-Nahentartung des Minimal-Even-Eigenvektors) ist das Bild qualitativ anders: $A_\lambda(0{,}4)\approx 1{,}004$ ist flach über $\lambda=3,5,7,9,11,13,15$ und $\sup|\hat\xi_\lambda|_{\mathbb R}\approx 0{,}88$ konstant. Die früheren dps=30-Messungen sind retrahiert; die methodische Lehre (dps-Bedarf skaliert linear mit $\lambda$ bei Cluster-nahentarteten Eigenvektoren) ist selbst ein substanzielles Ergebnis. Begleitend: EvenDom \texttt{K0\_NORMAL\_FAMILY\_REFRAME} \S E.18.
\end{itemize}

\begin{remark}[Gemeinsames Muster]
Diese fünf Ausschlüsse teilen ein strukturelles Muster: Der Minimal-Eigenvektor-Cluster von $QW_N$ ist die Quelle der technischen Schwierigkeit in jedem spektralen Versuch. Methoden, die eine uniforme Lücke voraussetzen (Kato--Temple, literales MS1), oder eine dünn-lokalisierte Eigenfunktion (naive Lokalisierung, Rand-Faktorisierung), oder ein generisches Funktionsklassen-Verhalten (dps-30 Streifenmessungen ohne Konvergenzverifikation) scheitern alle auf dieselbe Weise. Produktive Routen müssen daher \emph{mit} dem Cluster arbeiten (Waterline-/Cluster-Quotienten-Reformulierungen, wachstumsangepasste Normierungen), nicht dagegen.
\end{remark}


% =============================================================================
\section{Die vollständige Beweisarchitektur}
\label{sec:architecture}
% =============================================================================

Wir präsentieren die logische Kette vom bekannten Wissen zur Riemannschen Vermutung und identifizieren genau, wo jeder Teil passt und wo die Lücke bleibt.

\subsection{Die Reduktionskette}

\begin{center}
\renewcommand{\arraystretch}{1.3}
\resizebox{\textwidth}{!}{%
\begin{tabular}{c|p{7cm}|l|l}
\textbf{Schritt} & \textbf{Aussage} & \textbf{Status} & \textbf{Quelle} \\\hline
A1 & Connes' Theorem~6.1: gerade Dominanz $\Rightarrow$ Nullstellen von $\hat{\eta}$ reell & bewiesen & \cite{ConnesvS2025, Connes2026} \\
A2 & Hurwitz: gerade Dominanz für $\lambda_n \to \infty$ genügt & bewiesen & \cite{Connes2026} \\
A3 & Gerade Dominanz bei $33$ Werten, $\lambda \leq 1{,}300{,}000$ & numerische Evidenz / bedingte finite Brücke & Teil~II \\
A4 & Shift-Parity-Lemma: jede Primzahl bevorzugt gerade & bewiesen & Teil~II \\
A5 & Frontier-Primzahl-Mechanismus & bewiesen & Teil~II \\
A5b & Euler-Maclaurin: $\rho^{\mathrm{EM}}(13) < 0$ (intervallarithmetisch zertifiziert) & bewiesen & Teil~II \\
\textbf{A6} & \textbf{Kumulativer Schritt: $\sum_p W_p$ erhält die gerade Präferenz} & \textbf{variational / bedingt}${}^\dagger$ & Teil~II \\
A7 & Gerade Dominanz für $\lambda_n\to\infty$ & bedingt auf AVG / Odd-Lower-Bound & Teil~II \\
A8 & Alle Nullstellen auf $\Re(s) = 1/2$ (RH) & conditional reduction & A1+A2+A7
\end{tabular}%
}
\end{center}

\noindent${}^\dagger$\textbf{Status von Proposition~A6} (Teil~II, Proposition~A6): In v2.1 wurde ein robuster direkter Frontier-Dominanz-Beweis (Teil~II, Proposition~\emph{prop:direct-frontier}) hinzugefügt. Er verwendet einen gemeinsamen Rayleigh-Testvektor auf dem Frontier-Fenster $r\in[0.7,1]$, PNT-Partialsummation, den Satz von Mertens und finite-range CAP-Evidenz bis zur expliziten asymptotischen Schwelle. In der v2.3-Lesart beweist dies $\lambda_{\min}(W^+-W^-)<0$ asymptotisch, aber nicht automatisch $\lambda_1^+<\lambda_1^-$. Der Eigenwertordnungsschritt bleibt bedingt auf eine validierte untere Schranke für $\lambda_1^-$.

\subsection{Der kumulative Schritt (A6) --- bedingte Reduktion}

Der kumulative Schritt war die zentrale Herausforderung: zu zeigen, dass die Summe $\sum_{p \leq \lambda} W_p$ die individuelle gerade Präferenz (Shift-Parity-Lemma) als echte Eigenwertordnung beibehält. Teil~II trennt inzwischen die bewiesene Variationsdiagnostik von der noch offenen Eigenwertordnung:

\begin{enumerate}
  \item \textbf{Regime~1} ($\lambda \in [100, 1{,}300{,}000]$): $33$ computergestützte finite-range Diagnostiken mit Intervallarithmetik (50-Digit Präzision), theoremstark bedingt auf den Odd-Sektor-Tail-Lower-Bound.
  \item \textbf{Regime~2} (asymptotisch): Das M1$''$-Resolventen-Framework zeigt das korrekte Variationsvorzeichen. Drei Hilfsresultate kontrollieren die Approximationsqualität:
  \begin{itemize}
    \item \emph{Höhere-Moden-Abkling} (Lemma~B): Moden $j \geq 2$ tragen $O(D/L) = o(D)$ bei.
    \item \emph{Resolventen-Trunkierung} (Lemma~C): Neumann-Reihe konvergiert ($\|R_0 K\| < 1$ für $\lambda \geq 500$).
    \item \emph{PNT-Transfer} (Lemma~PNT): $|\rho - \rho^{\mathrm{EM}}| = O(L\,e^{-c\sqrt{L}})$.
  \end{itemize}
  \item \textbf{Überlappung}: Beide Regime treffen im Bereich $[442{,}413,\, 1{,}300{,}000]$ zusammen; die Restlücke ist eine Spektralrichtungs-Lücke, keine Abdeckungslücke.
\end{enumerate}

\smallskip
\noindent\textbf{v2.1-Ergänzung --- Direkte Frontier-Dominanz.} In v2.1 (Teil~II, Abschnitt zur direkten Frontier-Dominanz) wurde ein zweites Argument eingefügt. Seine robuste Fassung verwendet die Variationsungleichung $\lambda_{\min}(D_{\mathrm{total}})\leq e_1^\top D_{\mathrm{frontier}}e_1+\|D_{\mathrm{non}}\|$ mit dem gemeinsamen Vektor $e_1=(0,1,0)^\top$, statt einzelne Minimum-Eigenwerte zu addieren. Auf $r\in[0.7,1]$ erfüllt der skalare Frontier-Quotient $e_1^\top D_3(r)e_1\leq -0{,}4685$; PNT-Partialsummation liefert einen negativen Frontier-Beitrag der Ordnung $\sqrt{\lambda}$, während Mertens-artige Abschätzungen die Non-Frontier-Norm durch $O(\lambda^{0.35}+\log\lambda)$ begrenzen. Dieses Argument beweist die asymptotische Variationsaussage, aber nicht ohne weitere untere Schranke die Eigenwertordnung.

\subsection{Beweisstrategien für A6: Vergleichende Analyse}

Vier Kandidatenstrategien für den kumulativen Schritt wurden während der Forschung identifiziert:

\begin{center}
\footnotesize
\renewcommand{\arraystretch}{1.25}
\resizebox{\textwidth}{!}{%
\begin{tabular}{c|p{2.5cm}|p{2.8cm}|p{2.8cm}|l}
 & \textbf{Ziel} & \textbf{Stärke} & \textbf{Engpass} & \textbf{Entscheidung} \\\hline
\textbf{A} & Monotone EW-Ordnung via totale Positivität ($\mathrm{PF}_\infty$) & Wenn gültig: A6 automatisch (ordnungserhaltende Summen) & Fourier-Multiplikator $M(\xi) = -2\,\Re[\zeta'/\zeta(\tfrac{1}{2}+i\xi)]$ ist auf $\sim\!49\%$ seiner Domäne negativ; $\mathrm{PF}_\infty$ scheitert als absolute Eigenschaft & Ausgeschlossen (Thm.\ NE-A) \\[4pt]
\textbf{B} & Sturm-Liouville-Struktur via kommutierendem Operator & Wenn gefunden: Grundzustands-Rigidität macht A6 mechanisch & SVD über 13 $r$-Samples zeigt, dass der universelle Kommutator-Kern trivial ist ($T = c\cdot I$) für $N$ bis 15; $\sigma_2$ bleibt unabhängig von $N$ nach unten beschränkt & Ausgeschlossen für $N \leq 15$ (Thm.\ NE-B); voller Raum offen \\[4pt]
\textbf{C} & Stabilität des Grundzustands als prolate Störung & Connes' natürlicher Weg & Quantitative Störungskontrolle & \textbf{Backup} \\[4pt]
\textbf{D} & Globale Schranke via $\mathrm{Tr}(M^k)$ & Spuren sind additiv & Momente $\to$ Eigenwerte nicht gratis & \textbf{Backup} \\[4pt]\hline
\textbf{M1$''$} & Resolvent-gedämpfte Energie + PNT-Transfer & Nutzt direkt Shift Parity + 33 CAP & Braucht Lemma~B + C & \textbf{Gewählt}
\end{tabular}}
\end{center}

\noindent\textbf{Warum M1$''$ + PNT + CAP gewählt wurde.} Routen~A und~B wären jeweils ,,Jackpot''-Lösungen --- wenn sie funktionieren, ist A6 unmittelbar --- aber v1.4 zeigt, dass keine der beiden als geschlossener struktureller Mechanismus existiert. Route~A (totale Positivität) scheitert, weil der Fourier-Multiplikator des Primzahl-Shift-Operators, $M(\xi) = -2\,\Re[\zeta'/\zeta(\tfrac{1}{2}+i\xi)]$, nach der expliziten Formel auf etwa der Hälfte seiner Domäne negativ ist; die nicht-trivialen Nullstellen von $\zeta$ sind genau die Oszillationen von $M$. Route~B (kommutierender Operator) scheitert, weil eine SVD-Suche über ein dichtes $r$-Gitter zeigt, dass der einzige universell mit $D_N(r)$ kommutierende Operator die Identität ist, für alle $N \leq 15$ --- die Primzahl-Shifts auf verschiedenen Skalen haben strukturell verschiedene Eigenbasen und teilen keine gemeinsame algebraische Symmetrie.

Diese Ausschlüsse sind keine Beweisfehler, sondern strukturelle Entdeckungen: die gerade Dominanz gilt \emph{trotz} der Abwesenheit dieser beiden Eigenschaften, was klärt, dass der Mechanismus echt arithmetisch-statistisch und nicht algebraisch ist. Route~C (prolate Störung) ist die nächstliegende überlebende Alternative: das M1$''$-Framework ist strukturell ähnlich zur Behandlung der Primzahl-Beiträge als Störung. Route~D (Spur-Momente) liefert ein natürliches Backup, da Spuren additiv sind und die Eigenwert-Nicht-Additivität umgehen. Der gewählte Weg kombiniert die stärksten verfügbaren Werkzeuge:

\begin{itemize}
  \item \emph{Analytisches Rückgrat:} M1$''$ zeigt $\rho(\lambda) \to 0$ via Leading-Mode-Cancellation (Lemma~A) und PNT-Transfer.
  \item \emph{Numerische Brücke:} 33 finite-range Diagnostiken decken den endlichen Bereich bedingt ab, in dem die analytische Asymptotik noch nicht greift.
  \item \emph{Fehlerkontrolle:} Lemma~B (Höhere-Moden-Abkling) und Lemma~C (Neumann-Konvergenz) stellen sicher, dass die Resolventen-Approximation treu ist.
\end{itemize}

\noindent Routen~C und~D bleiben als unabhängige Verifikationspfade verfügbar. Insbesondere würde ein Spur-Momente-Beweis von Lemma~C die Abhängigkeit von der Neumann-Konvergenz-Bedingung $\|R_0 K\| < 1$ beseitigen und ein vollständig ,,weiches'' analytisches Argument für Regime~2 liefern.


% =============================================================================
\section{Die Lokalisierungstabelle}
\label{sec:localization}
% =============================================================================

Wir präsentieren zwei Perspektiven darauf, wie die Schwierigkeit der Riemannschen Vermutung verteilt ist.

\subsection{Li-Koeffizienten-Perspektive}

Aus der Bombieri-Lagarias-Zerlegung und der thermodynamischen Analyse:

\begin{center}
\renewcommand{\arraystretch}{1.2}
\resizebox{\textwidth}{!}{%
\begin{tabular}{c|l|l|l}
\textbf{Schritt} & \textbf{Aussage} & \textbf{Status} & \textbf{Methode} \\\hline
S1 & Definition von $\xi(s)$, Funktionalgleichung & bewiesen & Definition \\
S2 & Li-Koeffizienten via Bombieri-Lagarias & bewiesen & \cite{BombieriLagarias1999} \\
S3 & Zerlegung $\lambda_n = \lambda_n^{(\infty)} + \lambda_n^{(\mathrm{fin})}$ & bewiesen & Teil~I \\
S4 & Archimedische Dominanz: $\lambda_n^{(\infty)} \sim \frac{n}{2}\log n$ & bewiesen & Stirling \\
\textbf{S5} & \textbf{$\lambda_n \geq 0$ für alle $n \geq 1$} & \textbf{OFFEN}${}^{\ast}$ & \textbf{$\Longleftrightarrow$ RH} \\
S6 & Alle Nullstellen auf $\Re(s) = 1/2$ & $\Leftrightarrow$ S5 & Li (1997) \\
S7 & Primzahlzählung: $\pi(x) = \mathrm{Li}(x) + O(\sqrt{x}\log x)$ & $\Leftarrow$ S6 & von Koch
\end{tabular}%
}
\end{center}

{\footnotesize\noindent${}^{\ast}$ S5 ist äquivalent zu RH (Li, 1997). Der vorliegende Weg der geraden Dominanz liefert nur eine bedingte Reduktion bis zur Odd-Sektor-Lower-Bound-/AVG-Schließung; S5 ist hier daher nicht geschlossen.\par}

\subsection{Gerade-Dominanz-Perspektive}

Aus dem Connes-Programm und Teil~II:

\begin{center}
\renewcommand{\arraystretch}{1.2}
\resizebox{\textwidth}{!}{%
\begin{tabular}{c|l|l|l}
\textbf{Schritt} & \textbf{Aussage} & \textbf{Status} & \textbf{Methode} \\\hline
E1 & Explizite Weil-Formel & bewiesen & Weil (1952) \\
E2 & Quadratische Weil-Form $\mathrm{QW}_\lambda$ & bewiesen & Connes (2026) \\
E3 & Gerade/ungerade Zerlegung von $\mathrm{QW}_\lambda$ & bewiesen & Teil~II \\
E4 & Shift-Parity-Lemma: jede Primzahl bevorzugt gerade & bewiesen & Teil~II \\
E5 & Finite-range Gap-Diagnostik bei $33$ Werten, $\lambda \leq 1{,}300{,}000$ & Evidenz / bedingt & Teil~II (Diagnostik) \\
E5b & Euler-Maclaurin: $\rho^{\mathrm{EM}}(13) < 0$ (intervallarithmetisch zertifiziert) & bewiesen & Teil~II \\
\textbf{E6} & \textbf{Kumulativer Schritt: gerade Dominanz für alle $\lambda$} & \textbf{offen / bedingt} (= A6) & \textbf{Odd-Lower-Bound / AVG} \\
E7 & Connes' Theorem~6.1 + Hurwitz-Brücke $\Rightarrow$ RH & bedingt auf E6 & \cite{Connes2026}
\end{tabular}%
}
\end{center}

\noindent Die beiden Perspektiven sind verwandt, aber \emph{nicht äquivalent}: E6 (kumulative gerade Dominanz) würde S5 (und damit RH) via Connes' Theorem~6.1 und den Satz von Hurwitz \emph{implizieren}. Die Umkehrung ist jedoch unbekannt --- RH könnte aus Gründen gelten, die nichts mit gerader Dominanz zu tun haben. E6 ist eine \emph{hinreichende} Bedingung, keine Charakterisierung. Das Shift-Parity-Lemma überbrückt die Perspektiven, indem es zeigt, dass die Struktur der individuellen Primzahlen, die E6 zugrunde liegt, algebraisch kontrolliert ist.


% =============================================================================
\needspace{8cm}
\section{Systematisch geprüfte Alternativen}
\label{sec:explored-alternatives}
% =============================================================================

Während der Entwicklung des Beweises wurden elf alternative Strategien (BI-1 bis BI-11) systematisch exploriert. Diese erschöpfende Suche ist selbst ein Meta-Resultat: \emph{der Lösungsraum ist kartiert, und der gewählte Weg ist der einzig konsistente.}

\begin{center}
\footnotesize
\renewcommand{\arraystretch}{1.2}
\resizebox{\textwidth}{!}{%
\begin{tabular}{l|p{4.5cm}|l|l}
\textbf{Ansatz} & \textbf{Kernidee} & \textbf{Ergebnis} & \textbf{Status} \\\hline
BI-1: $\det_2$ & Raum erweitern (Hilbert--Schmidt) & Kein Positivitätsersatz & Widerlegt \\
BI-2: RG-Fixpunkt & Weil-Positivität unter Skalierung & Nicht formalisierbar & Verworfen \\
BI-3: Universeller Term & Physikalische Plausibilität $\Rightarrow$ RH & Zu abstrakt & Verworfen \\
BI-5: Lifting-Programme & Weil $\to$ Grothendieck $\to$ RMT & Offen außer $\mathbb{F}_q$ & Nicht nutzbar \\
BI-6: Implikationsring & Li $\leftrightarrow$ Weil $\leftrightarrow$ NB $\leftrightarrow$ BD & Ring existiert nicht & Widerlegt \\
BI-7: Pöschl--Teller & Gamma-Analogie zu $\xi$ & Keine Stabilität & Negativ \\
\textbf{BI-8: Stabilität zweiter Ordnung} & \textbf{RH als Variationsproblem} & \textbf{Architektur korrekt} & \textbf{Realisiert} \\
BI-9: Soliton-Modelle & Topologische Stabilität & Coulomb-Gas instabil & Widerlegt \\
\textbf{BI-10: Bausteine} & \textbf{No-Coordination, Shift-Struktur} & \textbf{Erfolgreich wiederverwendet} & \textbf{Integriert} \\
\textbf{BI-11: Connes 2026} & \textbf{Even-Dominanz via $Q_W$} & \textbf{Durchbruch} & \textbf{Kernweg}
\end{tabular}}
\end{center}

\noindent Die Schlüsselerkenntnis aus dieser Exploration: BI-8 (Stabilität zweiter Ordnung) war strukturell korrekt --- Positivität auf einem eingeschränkten Unterraum genügt, und die Entscheidung fällt in zweiter Ordnung --- aber mechanistisch falsch: die ,,Entropie'' in RH ist nicht thermodynamisch (KL-Divergenz), sondern arithmetisch-geometrisch (trigonometrische Struktur der Primzahl-Shifts, No-Coordination-Lemma). Das Shift-Parity-Lemma ersetzt den Varianz-/KL-Term vollständig.


% =============================================================================
\needspace{10cm}
\section{Unabhängige Resultate}
\label{sec:independent-results}
% =============================================================================

Mehrere Ergebnisse dieses Programms sind unabhängig vom Hauptbeweis von Interesse:

\begin{center}
\footnotesize
\renewcommand{\arraystretch}{1.2}
\resizebox{\textwidth}{!}{%
\begin{tabular}{p{4.2cm}|p{3.2cm}|p{4.2cm}|l}
\textbf{Resultat} & \textbf{Kontext} & \textbf{Bedeutung} & \textbf{Status} \\\hline
\multicolumn{4}{l}{\emph{Positive Resultate (neue Mathematik)}} \\[2pt]
Shift-Parity-Lemma & Arithm. Trigonometrie & Jede Primzahl bevorzugt individuell gerade Eigenfunktionen & Bewiesen \\
No-Coordination-Lemma & Multipl. Unabhängigkeit & Arithmetischer Ersatz für den Entropie-Term & Bewiesen \\
Leading-Mode-Cancellation ($c = 2+\sqrt{2}$) & Resolvent-Analyse & Universelle paarweise Cancellation-Konstante & Bewiesen \\
OS-Reflexionspositivität & Spektraltheorie & Bochner + kontraktive Halbgruppe & Bewiesen als struktureller Input \\
Euler--Maclaurin: $\rho^{\mathrm{EM}} < 0$ & PNT-Asymptotik & Globaler Stabilitätsanker & Bewiesen \\
Hellmann--Feynman: $\partial\Delta/\partial L > 0$ & Sensitivitätsanalyse & Lückenvertiefung durch neue Primzahlen, nicht durch~$L$ & Bewiesen \\
Zerlegung: Primzahlen treiben die variationsmäßige gerade Präferenz & Spektralzerlegung & Archimedischer Kern ist paritätsneutral & Variational / bedingt \\
$33$ finite-range Diagnostiken & Intervallarithmetik & Bedingte Brücke bis zur Odd-Lower-Bound-Validierung & Evidenz / bedingt \\[4pt]
\multicolumn{4}{l}{\emph{Negative Resultate (Obstruktionen und Unmöglichkeiten)}} \\[2pt]
Obstruktionen K1--K4 & Sackgassen-Analyse & Kartiert den Raum unmöglicher Strategien & Dokumentiert \\
Convexity Trap & Variationsmethoden & KL-Positivität $\not\Rightarrow$ Li-Positivität & Bewiesen \\
Weyl-Law-Obstruktion & Operatorentheorie & Kein naiver Hilbert--P\'olya-Operator existiert & Bewiesen \\
Mellin-Abkling-Barriere & Funktionalanalysis & Li-Polynome können Weil-Klasse nicht approximieren & Bewiesen \\
Implikationsring-Nichtexistenz & Kriterienvergleich & Li $\to$ Weil direkte Implikation unbekannt & Dokumentiert \\
Coulomb-Gas-Instabilität & Statistische Mechanik & Reine Positionsmodelle instabil abseits der Linie & Bewiesen \\
Kein absolutes $\mathrm{PF}_\infty$ für $A_\lambda$ & Primzahl-Shift-Multiplikator & $M(\xi) = -2\,\Re[\zeta'/\zeta(\tfrac{1}{2}+i\xi)]$ auf $\sim\!49\%$ der Domäne negativ & Bewiesen (v1.4) \\
Kein universeller kommutierender Operator & Sturm-Liouville-Suche & Kern von $[T, D_N(r)]$ ist $\{c \cdot I\}$ für alle getesteten $N$ & Bewiesen (v1.4) \\[4pt]
\multicolumn{4}{l}{\emph{Strukturelle Erkenntnisse (Meta-Resultate)}} \\[2pt]
Invertierte-Hesse-Struktur & YM/TU/RH-Vergleich & $\partial^2\lambda_n/\partial\delta^2 < 0$: RH ist ein Maximum & Bestätigt \\
Interdisziplinärer Katalog & Regularisierung & $11$ Felder teilen die Spurklassen-Obstruktion & Dokumentiert
\end{tabular}}
\end{center}

\noindent Diese $17$ Resultate zeigen, dass das Programm substanzielle Beiträge unabhängig vom Hauptbeweis hervorgebracht hat. Insbesondere sind das Shift-Parity-Lemma, die Obstruktionsanalyse (K1--K4) und die Mellin-Abkling-Barriere unabhängig publizierbar. Die negativen Resultate sind wohl die wertvollsten: sie schützen zukünftige Forscher davor, Sackgassen-Strategien zu wiederholen.


% =============================================================================
\section{Verbleibende offene Probleme}
\label{sec:open-problems}
% =============================================================================

Da OP1 nach v2.3 nur bedingt auf Odd-Lower-Bound/AVG reduziert ist, formulieren wir die verbleibenden offenen Probleme:

\begin{remark}[OP1: Kumulative gerade Dominanz --- bedingt auf Odd-Lower-Bound (v2.3)]
\label{rem:OP1-resolved}
Die kumulative gerade Dominanz (A6) ist nach v2.3 nicht als Eigenwertordnung geschlossen. Teil~II liefert zwei starke Bausteine: die v2.0-Drei-Regime-Diagnostik und den v2.1-direkten Frontier-Variationsbeweis (Teil~II, Proposition~\emph{prop:direct-frontier}). Das v2.0-Drei-Regime-Argument (siehe \Cref{sec:architecture}) kombiniert Regime~1 ($33$ finite-range Gap-Diagnostiken, bedingt auf den Odd-Sektor-Tail-Lower-Bound) und Regime~2 (M1$''$-Resolventen-Framework). Die Spektrallückenkonstante ist analytisch motiviert: $c_{\mathrm{gap}} \geq 4\pi^2/L \geq 2{,}8$ für $L \leq 14$ (Teil~II, Korollar nach Lemma~B, Schritt~3). Der v2.1-direkte Frontier-Dominanz-Beweis verwendet einen gemeinsamen Rayleigh-Testvektor auf dem Frontier-Fenster, PNT-Partialsummation, den Satz von Mertens und finite-range CAP-Evidenz bis zur expliziten asymptotischen Schwelle; dieses Argument beweist das Variationsvorzeichen, lässt den Schluss auf $\lambda_1^+<\lambda_1^-$ aber offen.
\end{remark}

\begin{remark}[OP2: Einfachheit von $\lambda_1^+$ --- Gelöst]
\label{rem:OP2-resolved}
Der kleinste Eigenwert $\lambda_1^+$ von $\mathrm{QW}_\lambda^+$ ist als einfach zertifiziert für alle $33$ Zertifikatwerte $\lambda \in [100,\, 1{,}300{,}000]$ mittels Intervallarithmetik auf dem $4 \times 4$ geraden Galerkin-Block (Teil~II, Lemma~Einfachheit). Die Intra-Even-Lücke $\lambda_2^+ - \lambda_1^+$ wächst monoton von $8{,}69$ ($\lambda = 100$) bis $\geq 731$ ($\lambda = 320{,}000$), skalierend als $\sim\!\sqrt{\lambda}$. Für Regime~2 ($\lambda \geq 442{,}413$) folgt die analytische Lücke aus der Galerkin-Diagonal-Asymptotik: $W_{11}^+$ ist der eindeutig negativste Diagonaleintrag. Sobald die Odd-Sektor-Lower-Bound-Schließung die Eigenwertordnung $\lambda_1^+ < \lambda_1^-$ liefert, zertifiziert dies, dass der globale Grundzustand von $A_\lambda$ einfach und gerade ist, wie von Connes' Theorem~6.1 gefordert.
\end{remark}

\begin{conjecture}[OP3: Analytischer Beweis der Indexstarre]
\label{conj:OP3}
Der negative Trägheitsindex von $\Ksym$ ist genau~$2$ für alle hinreichend großen $N$ und $\sigma \in (1, \sigma_{\max})$.
\end{conjecture}

\begin{conjecture}[OP4: Nuklearer Transferoperator für $\xi$]
\label{conj:OP4}
Man konstruiere einen separablen Hilbert-Raum $\mathcal{V}$ und eine nukleare Familie $s \mapsto \mathcal{L}_s$ mit $\xi(s)/\xi(0) = \det(I - \mathcal{L}_s)$, Spektralradius $< 1$ für $\Re(s) > 1/2$ und Eigenwert $1$ an den Nullstellen. Dies verbindet sich mit Deningers blättrigen Räumen und dem Hilbert-P\'olya-Programm.
\end{conjecture}

\begin{conjecture}[OP5: Nevanlinna-Eigenschaft von $\log\det_2$]
\label{conj:OP5}
Ist $\log\det_2(I - zR)$ eine Nevanlinna-Funktion (positiver Imaginärteil in der oberen Halbebene) genau dann, wenn alle Nullstellen auf der reellen Achse liegen? Dies wäre ein neues RH-äquivalentes Kriterium im $\det_2$-Rahmen (Teil~I, \S6).
\end{conjecture}

\noindent OP1 bleibt damit der zentrale offene Eigenwertordnungsschritt. OP2 ist finite-range bzw. strukturell diagnostiziert, aber für die asymptotische RH-Reduktion ebenfalls an die Odd-Lower-Bound-/AVG-Schließung gekoppelt. OP3--OP5 sind strukturelle Fragen, die alternative Ansätze oder tieferes Verständnis liefern könnten.


% =============================================================================
\section{Verbindungen zur weiteren Landschaft}
\label{sec:connections}
% =============================================================================

\subsection{Connes' Programm}

Die Arbeit in Teil~II fügt sich direkt in Connes' spektralen Ansatz zur Riemannschen Vermutung ein \cite{Connes2026}. Das Shift-Parity-Lemma liefert analytische Evidenz für die Hypothese der geraden Dominanz, die Connes in Theorem~6.1 voraussetzt. Die computergestützten Diagnostiken an $33$ Werten von $\lambda$ bis $1{,}300{,}000$ bieten starke finite-range Evidenz über mehr als $4$ Größenordnungen, bedingt auf die Odd-Sektor-Lower-Bound-Validierung. Der Resolventen-M1$''$-Rahmen, das Leading-Mode-Cancellation-Lemma und das PNT-Transfer-Lemma etablieren das asymptotische Variationsvorzeichen; die verbleibende Herausforderung ist die Eigenwertordnungs-/Lower-Bound-Schließung.

\subsection{Nyman-Beurling-Baez-Duarte}

Das Li-Kriterium, das Weil-Positivitätskriterium, das Nyman-Beurling-Kriterium und das Baez-Duarte-Kriterium sind alle äquivalent zur Riemannschen Vermutung. Direkte Implikationen zwischen ihnen (ohne den Umweg über RH) sind jedoch nur in zwei Fällen bekannt: Weil $\Rightarrow$ Li (Bombieri-Lagarias, 1999) und Nyman-Beurling $\Rightarrow$ Baez-Duarte (trivial). Ein geschlossener Implikationsring würde einen Beweis darstellen.

\subsection{Das Spurklassen-Hindernis als universelles Muster}

Das Hindernis K4 (Spurklassen-Versagen bei $\Re(s) = 1/2$) ist nicht einzigartig für die Riemannsche Zetafunktion. Die gleiche Dichtelücke zwischen Primzahlen ($\pi(x) \sim x/\log x$) und hyperbolischen Geodäten ($\#\{\gamma: l_\gamma \leq T\} \sim e^T/T$) erscheint in:
\begin{itemize}
  \item \textbf{QFT:} UV-Divergenzen, regularisiert durch Cutoff/Renormierung.
  \item \textbf{Statistische Mechanik:} Thermodynamischer Limes erfordert endliche-Volumen-Approximation.
  \item \textbf{NCG:} Connes' nichtkommutative Geometrie verwendet Zeta-Funktions-Regularisierung von Spektraltripeln.
\end{itemize}
In jedem Fall liefert ein externer Parameter (Temperatur, Cutoff, Skala) die Regularisierung. Die Zahlentheorie hat keinen solchen Parameter --- nur die Gammafunktion, die die Divergenz formal behebt, aber operatortheoretisch nicht realisiert wurde.

\subsection{Selberg-Zeta und Yang--Mills}

Die Selberg-Zetafunktion $Z_\Gamma(s)$ für eine kompakte hyperbolische Fläche $\Gamma \backslash \mathbb{H}$ hat dieselbe Euler-Produkt-Struktur wie Riemanns $\zeta(s)$, aber mit ``Primzahlen'' ersetzt durch primitive Geodäten. Das Spurklassen-Hindernis K4 greift \emph{nicht}: Geodätenlägen wachsen exponentiell ($\#\{\gamma: l_\gamma \leq T\} \sim e^T/T$), was Spurklassen-Konvergenz für alle $\Re(s) > 1/2$ sicherstellt. Die Selberg-Zeta-RH ist für kompakte Flächen bewiesen (Selberg, 1956).

Dieser Kontrast beleuchtet die arithmetische Schwierigkeit: Die Primzahlzählung der Riemannschen Zetafunktion ($\pi(x) \sim x/\log x$) erzeugt genau die ``falsche'' Dichte für Spurklassen-Argumente. Der Gerade-Dominanz-Ansatz umgeht dieses Hindernis, indem er mit der Weil-Quadratischen-Form (endlichdimensionale Galerkin-Blöcke) statt mit unendlichdimensionalen Transfer-Operatoren arbeitet.

Eine Yang--Mills-Analogie existiert: Kirks nullstellenfreier Streifen für Selberg-Zeta (via Ruelles Theorem) spiegelt die nullstellenfreie Region in der Primzahltheorie wider. Die Spektrallücke im Laplace-Operator auf $\Gamma \backslash \mathbb{H}$ spielt die Rolle der Gerade-Ungerade-Lücke in der Weil-Form. Ob der Shift-Parity-Mechanismus ein geometrisches Gegenstück für Selberg-Zeta hat, ist eine offene Frage von eigenständigem Interesse.

% added 2026-05-23: Route B / I-1 reframe / E.18 empirics
\subsection{Parallelprogramm: Connes--vS-Fehlerkollaps und I-1-Normalfamilien-Reframe}
\label{sec:parallel-route-b}

Unabhängig von der Asymptotischen Variationsgap-Vermutung (A7) liefert die Connes--van-Suijlekom-Fehlerkollaps-Route (Route~B) eine weitere Reduktion auf die Riemannsche Vermutung.
Ausgangspunkt ist dieselbe Aussage wie in Route~A: Connes--vS-Theorem~5.6 (= Schritt A1) liefert die finite-$N$-Schranke
$\mathrm{err}_\lambda \leq C e^{-c\lambda}$ für jede individuelle Riemann-Nullstelle, numerisch verifiziert auf $\lesssim 10^{-25\lambda}$ für $\geq 15$ Nullstellen.
Route~B reformuliert dann die Limitesbedingung: es gilt
$\mathrm{err} \to 0 \Longleftrightarrow \text{(i)} + \text{(ii)}$
rigoros, wobei (i) Schritt~A1 ist und (ii) Connes' \S6.6(b)/MS2-Bedingung bezeichnet (eine uniforme Streifenschranke an die Familie $\{\hat\xi_\lambda\}$ kanonischer Kriteriumsfunktionen).
Da $\text{(ii)} \Rightarrow \mathrm{RH}$ ebenfalls rigoros ist, reduziert Route~B auf die einzige offene Bedingung~(ii) \cite{Geiger2026evendom}.

Der I-1-Normalfamilien-Reframe (\S E.1--E.18 der Begleit-Beweisnotizen) zerlegt~(ii) über ein Montel--Vitali-Kompaktheitsargument in zwei Teilbedingungen:
(ii-a)~eine uniforme wachstumsangepasste Streifenschranke an $\{\hat\xi_\lambda\}$ und
(ii-c)~eine Limesidentifikationsbedingung.
Der algebraische Kernschritt (\S E.14) zeigt, dass (ii-c) rigoros in~(ii-a) absorbiert wird: die Parität von $\mathrm{QW}_N$ (reell-symmetrisch, Gerade-Paritäts-Projektion $\Rightarrow \hat\xi_\lambda$ ist gerade und reell-auf-reell) erzwingt zusammen mit einem Hadamard-Ordnungsargument (unter der Annahme~(Z), dass alle Nullstellen einfach sind) das identische Verschwinden des Bohr--Carathéodory-Phasenfaktors.
Das kanonische Kriteriumsobjekt ist definitiv als
$\hat\xi(z) = \tfrac{2}{\sqrt{L}}\sin\!\bigl(\tfrac{zL}{2}\bigr)F(z)$
festgelegt, wobei $F(z) = \sum_{n \leq N} \Lambda(n)/\sqrt{n}\cdot\cos(z\log n/2)$
(zentriert, gerade, reell-auf-reell, ganz; Abweichung von numerischen Daten
$\leq 1{,}9 \times 10^{-52}$ bei $\lambda = 5$, $7$, $\mathrm{dps}=50$, \S E.15).
Route~B reduziert sich damit vollständig auf die offene Bedingung~(ii-a): eine rigorose uniforme Schranke aus $\mathrm{QW}_N$-Koerzitivität.

Konvergierte empirische Daten für das Sup-Observable
$A_\lambda(\delta) := \sup_{x\in\mathbb{R}} |\hat\xi_\lambda(x+i\delta)|/|\hat\xi_\lambda(x)|$
und das $L^2$-Observable $B_\lambda(\delta)$ bei $\delta = 0{,}4$, erhoben auf einem dedizierten Rechenknoten
(Mac~Studio, $\mathrm{dps}_{\min} \approx 30\lambda$--$50\lambda$
zur Kontrolle von Cluster-Nahentartung; zwei verschiedene dps-Werte bit-identisch bei $\lambda = 11, 13, 15$),
sind in Tabelle~\ref{tab:route-b-empirics} an sieben konvergierten Werten
$\lambda \in \{3, 5, 7, 9, 11, 13, 15\}$ zusammengefasst.

\begin{table}[h]
\centering
\small
\begin{tabular}{r r r r r}
\hline
$\lambda$ & $\sup|\hat\xi_\lambda|_{\mathbb{R}}$ & $A_\lambda(0{,}4)$ & $B_\lambda(0{,}4)$ & $B_\lambda(0{,}4)/\lambda^{0{,}4}$ \\
\hline
 3 & 0{,}858 & 1{,}0033  & 1{,}151 & 0{,}74 \\
 5 & 0{,}879 & 1{,}0037  & 1{,}340 & 0{,}70 \\
 7 & 0{,}886 & 1{,}0038  & 1{,}505 & 0{,}69 \\
 9 & 0{,}890 & 1{,}0039  & 1{,}649 & 0{,}68 \\
11 & 0{,}890 & 1{,}0039  & 1{,}778 & 0{,}68 \\
13 & 0{,}891 & 1{,}0039  & 1{,}895 & 0{,}68 \\
15 & 0{,}892 & 1{,}00394 & 2{,}003 & 0{,}68 \\
\hline
\end{tabular}
\caption{Konvergierte Empirie der Route-B-Observablen bei $\delta = 0{,}4$
(Begleit-Beweisnotizen, \S E.18). Das Sup-Observable $A_\lambda(0{,}4)$ ist
im Wesentlichen flach bei $\approx 1{,}004$, im scharfen Kontrast zur
generischen Phragmén--Lindelöf-Vorhersage $\sim \lambda^\delta$. Das
$B_\lambda$-Verhältnis fällt von $0{,}74$ bei $\lambda = 3$ auf ein stabiles
Plateau bei $0{,}68$ für $\lambda \geq 9$.}
\label{tab:route-b-empirics}
\end{table}

Die beobachtete Beinahe-Flachheit von $A_\lambda(0{,}4)$ zeigt, dass die Familie $\{\hat\xi_\lambda\}$
stark subgenerisch ist, was~(ii-a) empirisch plausibel macht.
Das $L^2$-Observable wächst als $B_\lambda(0{,}4) \approx 0{,}68\,\lambda^{0{,}4}$ auf dem Plateau,
konsistent mit polynomialem subgenerischen Verhalten.
Frühere Läufe bei $\mathrm{dps}=30$ (\S E.16--E.17) wurden wegen unter-konvergierter Eigenvektoren retrahiert; die korrigierten Daten sind in \S E.18
der Begleit-Beweisnotizen dokumentiert \cite{Geiger2026evendom}.


% =============================================================================
\section{Einschätzung und Ausblick}
\label{sec:outlook}
% =============================================================================

\subsection{Was wir erreicht haben}

\begin{enumerate}
  \item Eine \textbf{vollständige Strukturanalyse} des eichtheoretischen Ansatzes zur Riemannschen Vermutung, mit neun unbedingten Ergebnissen (R1--R9) und vier definitiv ausgeschlossenen Wegen (K1--K4).

  \item Das \textbf{Shift-Parity-Lemma}, ein rein algebraisches Ergebnis über trigonometrische Shift-Matrizen, das die universelle gerade Präferenz einzelner Primzahlen etabliert.

  \item Eine \textbf{computergestützte finite-range Diagnostik} der geraden Dominanz an $33$ Werten ($\lambda = 100$ bis $1{,}300{,}000$), die starke Evidenz für Connes' spektrale Hypothese über mehr als $4$ Größenordnungen liefert und theoremstark auf die Odd-Sektor-Lower-Bound-Validierung angewiesen ist.

  \item Die \textbf{präzise Reduktion des kumulativen Schritts (A6/E6)} auf eine Variationsdiagnostik plus Odd-Sektor-Lower-Bound: (i) das v2.0-Drei-Regime-Argument kombiniert computergestützte finite-range Diagnostik ($\lambda \leq 1{,}300{,}000$), das M1$''$-Resolventen-Framework mit PNT-Transfer und drei Kontroll-Lemmata; (ii) das v2.1-Argument der direkten Frontier-Dominanz (Teil~II, Proposition~\emph{prop:direct-frontier}) verwendet einen gemeinsamen Rayleigh-Testvektor, PNT+Mertens-Partialsummation und finite-range CAP-Evidenz. Zusammen lokalisieren diese die verbleibende Lücke als Eigenwertordnungsproblem, nicht als fehlende Frontier- oder PNT-Skala.

  \item Die \textbf{Euler-Maclaurin-Proposition}, die $\rho^{\mathrm{EM}}(13) < 0$ etabliert (intervallarithmetisch zertifiziert) und damit das analytische Rückgrat für Regime~2 des Brückenarguments liefert.

  \item \textbf{Strukturelle Erkenntnisse} (Frontier-Primzahl-Mechanismus, Hellmann-Feynman-Analyse, Lücken-Asymptotik), die jeden zukünftigen Beweis einschränken und leiten.
\end{enumerate}

\subsection{Der Weg zum Abschluss der Beweislücke}
\label{sec:gap-closure-path}

Die Euler-Maclaurin-Proposition (Teil~II, Proposition~4.11) etabliert, dass die PNT-kontinuierliche Approximation $\rho^{\mathrm{EM}}(L) < 0$ für $L \geq 13$, d.\,h. für $\lambda \geq e^{13} \approx 442{,}413$, gilt (intervallarithmetisch zertifiziert). Dies bedeutet, dass im kontinuierlichen (Primzahlen-durch-Integrale-ersetzten) Limes das relevante Variationsvorzeichen für alle hinreichend großen $\lambda$ gilt. Um dies als Eigenwertordnung für die tatsächlichen Primsummen zu nutzen, werden zwei Komponenten benötigt:

\begin{enumerate}
  \item \textbf{Computergestützte Diagnostik ($\lambda \leq 1{,}300{,}000$):} Starke bedingte Brücke für $33$ Werte, bis der Odd-Sektor-Lower-Bound validiert ist (Teil~II, \S3).
  \item \textbf{PNT-Transfer-Lemma (Teil~II):} Durch Abel-Summation mit dem Chebyshev-Funktionsfehler $\theta(x) - x = O(x\exp(-c\sqrt{\log x}))$ erfüllt das tatsächliche Resolventen-Verhältnis $\rho(\lambda) = \rho^{\mathrm{EM}}(\lambda) + O(L\,e^{-c\sqrt{L}}) \to \rho^{\mathrm{EM}}(\lambda)$ für $\lambda \to \infty$. Da $\rho^{\mathrm{EM}} < 0$ für $L \geq 14$ (intervallarithmetisch zertifiziert), existiert $\lambda_0$ mit $\rho(\lambda) < 0$ für alle $\lambda \geq \lambda_0$.
\end{enumerate}

\noindent Der Status von M1$''$ (Resolventen-Subdominanz) ist somit: Es liefert das richtige asymptotische Variationsvorzeichen. In v2.1 wird diese Variationsroute zusätzlich durch den robusten direkten Frontier-Dominanz-Beweis (Teil~II, Proposition~\emph{prop:direct-frontier}) gestützt. Kombiniert mit der Höhere-Moden-Abklingschranke (Lemma~B), der Resolventen-Trunkierungskontrolle (Lemma~C) und den $33$ computergestützten finite-range Diagnostiken ergibt sich Proposition~A6 als conditional reduction: Gerade Dominanz als Eigenwertordnung bleibt bedingt auf die Asymptotische Variationsgap-Vermutung bzw. eine validierte untere Schranke für $\lambda_1^-$.

\subsection{Analytischer Status der Kontroll-Lemmata (v1.1)}

In v1.1 wurden Lemma~B und Lemma~C zu vollständig analytischen Argumenten aufgewertet:
\begin{enumerate}
  \item \textbf{Lemma~B (Parity-Split-Lücke):} Die Spektrallücke hat eine Zwei-Stufen-Struktur, abgeleitet aus der Laplace-Asymptotik der Auto-Overlap-Profile. Gerade-Index-Moden haben Lücken $\sim 2\sqrt{\lambda}$, ungerade-Index-Moden $\sim 4\pi^2(j^2-1)\sqrt{\lambda}/L^2$. Korollar: $c_{\mathrm{gap}} \geq 4\pi^2/L \geq 2{,}8$ für $L \leq 14$ (Teil~II, Lemma~4.10, Schritt~3).
  \item \textbf{Lemma~C (Sektor-Differenz-Kontrolle):} Die globale Bedingung $\|R_0 K\| < 1$ wurde durch die schwächere und direkt relevante Aussage ersetzt, dass die Tail-Beiträge zu $E_{\sin}$ und $E_{\cos}$ nahezu gleich sind (durch Paritätssymmetrie bei führender Laplace-Ordnung), sodass ihre Differenz $O(D/N_0^2) \ll D$ ist (Teil~II, Lemma~4.12).
\end{enumerate}
Keine numerischen Konstanten verbleiben im Beweis von Regime~2. Für präzise Verweise: Lemma~B, Schritt~3 (Teil~II, Lemma~4.10) leitet die Parity-Split-Lücke ab; Lemma~C (Teil~II, Lemma~4.12) ersetzt die globale Neumann-Bedingung durch eine Sektor-Differenz-Schranke von $O(1/N_0^2)$.

\subsection{Die verbleibende Distanz}

Die Reduktionskette A1--A8 (\Cref{sec:architecture}) ist als \emph{conditional reduction} vollständig lokalisiert. Der kumulative Schritt~A6, zuvor die einzige verbleibende Herausforderung, ist in Teil~II (Proposition~A6) nicht als Eigenwertordnung gelöst, sondern auf eine Variationsdiagnostik plus Odd-Sektor-Lower-Bound reduziert. Das Drei-Regime-Argument kombiniert computergestützte finite-range Diagnostik (Regime~1, $\lambda \leq 1{,}300{,}000$), das M1$''$-Resolventen-Framework mit PNT-Transfer (Regime~2) und drei Kontroll-Lemmata: Höhere-Moden-Abkling (Lemma~B), Resolventen-Trunkierung (Lemma~C) und PNT-Fehler-Transfer (Lemma~PNT). In v2.1 wird ein direkter Frontier-Dominanz-Beweis (Teil~II, Proposition~\emph{prop:direct-frontier}) hinzugefügt, der die Interpolations- und PNT-Transfer-Caveats für die Variationsaussage durch einen gemeinsamen Frontier-Rayleigh-Vektor umgeht.

\subsection{Philosophische Reflexion: Genese der Beweisarchitektur}

Die Beweisarchitektur entstand nicht aus einer linearen Deduktion, sondern aus einer systematischen Erkundung strukturell verwandter Stabilitätsprobleme. Der Ausgangspunkt (BI-8) war die Beobachtung, dass drei \emph{a priori} unzusammenhängende Kontexte --- Yang--Mills-Theorie, Turbulenztheorie und die Riemannsche Vermutung --- dieselbe formale Architektur teilen: in jedem Fall entscheidet die Positivität einer zweiten Variation über Stabilität oder Regularität. In Yang--Mills und der Turbulenz folgt diese Positivität aus einem Entropie- oder Varianz-Term. Im arithmetischen Kontext existiert kein solcher probabilistischer Freiheitsgrad.

Frühe numerische Experimente zeigten, dass der naive Transfer auf RH in \emph{invertierter Form} erscheint: die Li-Koeffizienten $\lambda_n$ haben ein \emph{Maximum} bei der kritischen Konfiguration ($\delta = 0$, d.\,h. RH), kein Minimum. Das korrekte Funktional ist daher $F_{\mathrm{RH}} = -\lambda_n$, mit Minimum bei RH. Diese Inversion wies auf einen tieferen strukturellen Mechanismus hin: Stabilität kann hier nicht aus Konvexität im klassischen (KL-Divergenz-)Sinne kommen.

Soliton- und Coulomb-Gas-Modelle (BI-9) wurden als topologische Alternativen getestet. Alle scheiterten: Das Coulomb-Gas zeigte universell $\partial^2 E/\partial\delta^2 < 0$, was bedeutet, dass logarithmische Abstoßung die On-Line-Konfiguration \emph{destabilisiert}. Die entscheidende Erkenntnis war, dass Modelle, die nur die \emph{Positionen} der Nullstellen erfassen, ohne deren analytische Struktur (Ganzheit, Funktionalgleichung, Euler-Produkt), keine Stabilität liefern können. Lokale Potentiale sind unzureichend.

Eine kritische Überprüfung früherer Arbeiten (BI-10) identifizierte drei dauerhafte Bausteine: die Weyl-Law-Obstruktion (kein naiver Hilbert--P\'olya-Operator), die Convexity Trap (KL-Positivität impliziert nicht Li-Positivität) und das No-Coordination-Lemma (multiplikative Unabhängigkeit der Primzahlen verhindert kohärente Phasensynchronisation). Letzteres erwies sich als der arithmetische Ersatz für den fehlenden Entropie-Term.

Der Durchbruch kam mit Connes' 2026-Übersichtsarbeit \cite{Connes2026} (BI-11), die einen völlig anderen Rahmen lieferte. Anstelle von Dichteargumenten oder globaler Positivität wird die Weil-Form als endlicher Operator auf einem Paley--Wiener-Raum behandelt. Connes' Theorem~6.1 reduziert RH auf eine Gerade-Dominanz-Bedingung, und ein Hurwitz-Argument zeigt, dass diese Bedingung nur entlang einer Folge $\lambda_n \to \infty$ benötigt wird. Dies realisierte präzise die von BI-8 antizipierte ``Positivität auf einem eingeschränkten Unterraum''.

Die Beweisarchitektur von Teil~II ist die direkte Umsetzung dieser Erkenntnisse. Das Shift-Parity-Lemma formalisiert die konstruktive Interferenz gerader Moden und destruktive Interferenz ungerader Moden unter Primzahl-Shifts. Die Resolventenanalyse (M1$''$) und die Hilfslemmata über Modenabkling und Trunkierung zeigen, dass diese lokale gerade Präferenz das richtige Variationsvorzeichen unter Summation hat. Proposition~A6 lokalisiert die Brücke zur vollständigen Weil-Form, schließt aber den Odd-Sektor-Eigenwertordnungsschritt noch nicht.

Rückblickend sagte BI-8 die Architektur korrekt voraus, identifizierte aber den Mechanismus falsch. Die ``Entropie'' der Riemannschen Vermutung ist nicht thermodynamisch, sondern geometrisch-arithmetisch: sie entsteht aus der trigonometrischen Struktur der Primzahl-Shifts und der Nicht-Koordinierbarkeit ihrer Phasen. Die vorliegende Arbeit kann daher als Destillation dieser Entwicklung verstanden werden, in der alle heuristischen Elemente durch formale, verifizierbare Argumente ersetzt wurden.

Die v1.4-Ausschlüsse der Routen~A und~B schärfen dieses Bild weiter. Man hätte erwarten können, dass die Primzahlen, die durch $A_\lambda$ wirken, entweder einem absoluten Glättungsprinzip (totale Positivität) oder einem verborgenen Symmetrieprinzip (kommutierender Operator) gehorchen. Keines von beiden gilt: Der Fourier-Multiplikator $M(\xi)$ des Primzahl-Shift-Operators oszilliert durch Null, weil die nicht-trivialen Nullstellen von $\zeta$ seine Pole sind, und keine feste algebraische Struktur koordiniert die Shifts auf verschiedenen Skalen. Was übrig bleibt, ist das punktweise Shift-Parity-Lemma --- eine algebraische Identität trigonometrischer Überlappungen --- kombiniert mit der statistischen Konzentration der Primzahl-Gewichtung an der Frontier $p \sim \lambda$. Die Abwesenheit eines globalen Glättungs- oder Symmetrieprinzips ist keine zu füllende Lücke; sie ist die korrekte Struktur. Gerade Dominanz entsteht, weil die Primzahlen multiplikativ unabhängig sind, nicht trotz dieser Tatsache.

In v1.1 wurden beide verbliebenen numerischen Inputs analytisch ersetzt: Die Parity-Split-Lücke (Lemma~B) liefert $c_{\mathrm{gap}} \geq 4\pi^2/L$, und die Sektor-Differenz-Kontrolle (Lemma~C) ersetzt die globale Neumann-Bedingung. Keine numerischen Konstanten verbleiben im Beweis von Regime~2.


% =============================================================================
\section{Zusammenfassung}
\label{sec:summary}
% =============================================================================

\begin{center}
\renewcommand{\arraystretch}{1.3}
\resizebox{\textwidth}{!}{%
\begin{tabular}{l|c|l}
\textbf{Beitrag} & \textbf{Status} & \textbf{Arbeit} \\\hline
Thermodynamischer Formalismus für $\zeta$ & etabliert & I \\
9 unbedingte strukturelle Ergebnisse (R1--R9) & bewiesen & I \\
4 ausgeschlossene Wege (K1--K4) & bewiesen & I \\
Lokalisierung: gesamte RH $\equiv$ Schritt S5 & bewiesen & I \\
Shift-Parity-Lemma & bewiesen (analytisch) & II \\
Finite-range Gap-Diagnostik ($33$ Werte, $\lambda \leq 1{,}300{,}000$) & Evidenz / bedingt & II \\
Frontier-Primzahl-Mechanismus & bewiesen & II \\
Hellmann-Feynman: Primzahlzählung treibt Lücke & bewiesen & II \\
Leading-Mode-Cancellation ($c = 2+\sqrt{2}$) & bewiesen & II \\
Euler-Maclaurin: $\rho^{\mathrm{EM}} < 0$ für $\lambda \geq 1{,}2\mathrm{M}$ & bewiesen & II \\
Resolventen-M1$''$: $(E_{\sin}-E_{\cos})/D = O(1/L)$ & variational etabliert & II \\
Lücken-Asymptotik: $|\mathrm{cert\_gap}| \sim c\sqrt{\lambda}$ & numerisch & II \\
\textbf{M1$''$ (Resolventen-Subdominanz)} & \textbf{Variationsreduktion} & II \\
Höhere-Moden-Abkling (Lemma~B) & bewiesen & II \\
Resolventen-Trunkierung (Lemma~C) & bewiesen (analytisch) & II \\
PNT-Transfer-Fehler (Lemma~PNT) & bewiesen mit effektivem $C$ & II \\
\textbf{Kumulativer Schritt (A6)} & \textbf{variational / bedingt} & II \\
\textbf{Riemannsche Vermutung} & \textbf{conditional reduction} & I+II+III
\end{tabular}%
}
\end{center}

\noindent A6 beruht auf dem Drei-Regime-Brückenargument (Teil~II, Proposition~A6): computergestützte finite-range Diagnostik für $\lambda \leq 1{,}300{,}000$ (Regime~1, bedingt auf Odd-Sektor-Tail-Lower-Bound) und das analytische M1$''$/PNT-Argument für das asymptotische Variationsvorzeichen. In v1.1 wurde Lemma~B zu einer Parity-Split-Lückenschranke aufgewertet und Lemma~C als Sektor-Differenz-Kontrolle reformuliert. In v2.1 (Teil~II, Proposition~\emph{prop:direct-frontier}) wird ein robuster direkter Frontier-Variationsbeweis hinzugefügt: Die Frontier wird auf einem gemeinsamen Rayleigh-Vektor ausgewertet und das Komplement normbeschränkt. Beide alten Caveats des v2.0-Drei-Regime-Arguments entfallen für die Variationsaussage; der Beweis der Riemannschen Vermutung bleibt bedingt auf die Eigenwertordnungs-Schließung.

\bigskip
\noindent\textbf{Danksagung.} Berechnungen wurden auf einem Hetzner-CCX13-Cloud-Server (2~dedizierte vCPU, 8~GB RAM) durchgeführt. Die Intervallarithmetik-Bibliothek mpmath wurde von Fredrik Johansson entwickelt.


\section*{KI-Offenlegung}

Dieses Manuskript wurde in umfassender Zusammenarbeit mit KI-Sprachmodellen (Claude, Anthropic; Copilot, Microsoft/GitHub; Gemini, Google DeepMind) erstellt. Die KI trug zur konzeptionellen Exploration, analytischen Arbeit, Literaturrecherche, Texterstellung und kritischen Überprüfung bei. Der Autor, Lukas Geiger, hat die Forschungsrichtung konzipiert, seine Fachexpertise eingebracht und die finale redaktionelle Entscheidungsgewalt ausgeübt. Der Autor übernimmt die alleinige wissenschaftliche Verantwortung für den gesamten Inhalt einschließlich etwaiger Fehler.


\begin{thebibliography}{99}

\bibitem{Geiger2026partI}
L.~Geiger, \emph{Die Riemannsche Vermutung: Grundlagen, Hindernisse und Neuausrichtung}, 2026.

\bibitem{Geiger2026partII}
L.~Geiger, \emph{Die Riemannsche Vermutung: Gerade Dominanz der quadratischen Weil-Form}, 2026.

\bibitem{Connes2026}
A.~Connes, \emph{The Riemann Hypothesis: Past, Present and a Letter Through Time}, arXiv:2602.04022v1, 2026.

\bibitem{ConnesvS2025}
A.~Connes und W.~D. van Suijlekom, \emph{Quadratic Forms, Real Zeros and Echoes of the Spectral Action}, Commun.~Math.~Phys.~\textbf{406}, Article~312 (2025), DOI: 10.1007/s00220-025-05493-1.

\bibitem{Geiger2026evendom}
L.~Geiger, \emph{A conditional proof of the Riemann hypothesis via even dominance of the Weil quadratic form}, Zenodo, 2026, DOI: 10.5281/zenodo.20291994. % added 2026-05-23

\bibitem{BombieriLagarias1999}
E.~Bombieri und J.~C.~Lagarias, \emph{Complements to Li's criterion for the Riemann hypothesis}, J.~Number Theory~\textbf{77} (1999), 274--287.

\end{thebibliography}

\end{document}
